Neural thickets --- dense clusters of task-solving expert solutions reachable via random weight perturbations around pretrained models --- offer a remarkably cheap approach to ensemble construction. While this phenomenon has been demonstrated for large language models on reasoning tasks, its implications for adversarial robustness and out-of-distribution (OOD) detection remain unexplored. We transplant the neural thicket ensemble concept to the vision domain, constructing ensembles of up to 10 members by perturbing a pretrained ResNet-18 on CIFAR-10. We evaluate three perturbation strategies (isotropic Gaussian, orthogonal, and layer-scaled) under white-box PGD-20 and FGSM attacks, and assess OOD detection performance on SVHN using energy scores and mutual information. Our results reveal a fundamental limitation: vision models exhibit a sharp sensitivity cliff, where perturbation scales large enough to create meaningful diversity ($\sigma \geq 0.02$) catastrophically destroy accuracy (94\% $\to$ 28\%), while scales that preserve accuracy ($\sigma \leq 0.001$) produce near-identical ensemble members with $<$1\% disagreement. Consequently, thicket ensembles achieve PGD-20 accuracy of 12.5\%, indistinguishable from a single model (12.6\%), while adversarial training reaches 78.5\%. Neither orthogonal nor layer-scaled perturbations overcome this limitation. We conclude that the neural thicket phenomenon may be specific to heavily overparameterized models, and that weight perturbation alone is insufficient for adversarial robustness in vision.
